\subsection{CONTINUOUS MORPHISMS}

THEOREM 1. Let G and H be two Lie group germs over R or \(Q_{p}\) . Let f be a continuous morphism of G into H. Then f is analytic.

We give \(\mathrm{L}(\mathrm{G})\) and \(\mathrm{L}(\mathrm{H})\) norms which define their topologies and such that \(\| [x,y]\| \leqslant \| x\| \| y\|\) for all \(x,y\) . There exist an open ball V of centre 0 in \(\mathrm{L}(\mathrm{G})\) and an exponential mapping \(\phi\) of \(\mathbf{G}\) defined on \(\mathbf{V}\) such that: (1) \(\phi (\mathbf{V})\) is an open neighbourhood of \(e\) in G; (2) \(\phi\) is an isomorphism of the analytic manifold V onto the analytic manifold \(\phi (\mathbf{V})\) ; (3) \(\phi (nx) = \phi (x)^n\) for all \(x\in V\) and all \(n\in \mathbf{Z}\) such that \(nx\in V\) . We define similarly W and \(\psi\) for H. By shrinking V if necessary, it can be assumed that \(f(\phi (\mathrm{V}))\subset \psi (\mathrm{W})\) . Then \(g = \psi^{-1}\circ f\circ \phi\) is a continuous mapping of V into W.

We show that

\[
(x \in \mathrm{V}, \lambda \in \mathbf {Q} \text {   and   } \lambda x \in \mathrm{V}) \Rightarrow g (\lambda x) = \lambda g (x).\tag{1}
\]

It can be assumed that \(\lambda \neq 0\) . Let \(\lambda = \frac{r}{q}\) with \(r, q\) in \(\mathbf{Z} - \{0\}\) . Let \(y = \frac{r}{q} x\) .

If \(\mathbf{K} = \mathbf{R}\) , we write \(z = \frac{x}{q} = \frac{y}{r} \in \mathrm{V}\) . Then \(x = qz, y = rz\) , whence

\[
g (x) = \psi^ {- 1} (f (\phi (q z))) = \psi^ {- 1} (f (\phi (z) ^ {q})) = \psi^ {- 1} (f (\phi (z)) ^ {q}).
\]

We show that \(\psi^{-1}(f(\phi(z))^q) = q\psi^{-1}(f(\phi(z))) = qg(z)\) . It suffices to verify that, if \(u \in \psi(W)\) is such that \(u^q \in \psi(W)\) , then \(\psi^{-1}(u^q) = q\psi^{-1}(u)\) ; but if \(u = \psi(v)\) , \(u^q = \psi(v')\) , then \(v'/q \in W\) and \((\psi(v'/q))^q = u^q\) , hence \(\psi(v'/q) = u = \psi(v)\) and therefore \(v' = qv\) .

Similarly, \(g(y) = rg(z)\) , whence (1).

If \(\mathbf{K} = \mathbf{Q}_p\) , we write \(z = rx = qy \in \mathrm{V}\) , whence \(g(z) = rg(x) = qg(y)\) , whence again (1).

As \(\mathbf{Q}\) is dense in K, (1) implies that

\[
(x \in \mathrm{V}, \lambda \in \mathrm{K} \text {   and   } \lambda x \in \mathrm{V}) \Rightarrow g (\lambda x) = \lambda g (x).\tag{2}
\]

Let \(x \in L(G)\) and \(\lambda, \lambda'\) be elements of \(K^*\) such that \(\lambda x \in V, \lambda' x \in V\) . Then

\[
g (\lambda^ {\prime} x) = g \left(\frac {\lambda^ {\prime}}{\lambda} \lambda x\right) = \frac {\lambda^ {\prime}}{\lambda} g (\lambda x)
\]

by (2) and hence \(\frac{1}{\lambda} g(\lambda x) = \frac{1}{\lambda'} g(\lambda' x)\) . Thus an extension \(h\) of \(g\) to \(\mathrm{L}(G)\) is defined by writing \(h(x) = \frac{1}{\lambda} g(\lambda x)\) for all \(\lambda\) such that \(\lambda x \in \mathrm{V}\) . Clearly \(h\) is continuous. We show that

\[
(x \in \mathrm{L} (\mathrm{G}) \text {   and   } \lambda \in \mathrm{K}) \Rightarrow h (\lambda x) = \lambda h (x).\tag{3}
\]

Let \(\lambda' \in \mathbf{K}^*\) be such that \(\lambda' x \in V\) and \(\lambda' \lambda x \in V\) . Then

\[
h (\lambda x) = \frac {1}{\lambda^ {\prime}} g \left(\lambda^ {\prime} \lambda x\right) = \frac {1}{\lambda^ {\prime}} \lambda g \left(\lambda^ {\prime} x\right) = \lambda \frac {1}{\lambda^ {\prime}} g \left(\lambda^ {\prime} x\right) = \lambda h (x).
\]

Let \(x, y\) be in L(G). Then, by Proposition 4 of § 4, no. 3,

\[
\begin{array}{r l} h (x) + h (y) & = \lim _ {\lambda \in \mathbf {K} ^ {*},   \lambda \to 0} \lambda^ {- 1} \psi^ {- 1} (\psi (\lambda h (x)) \psi (\lambda h (y))) \\ & = \lim _ {\lambda \in \mathbf {K} ^ {*},   \lambda \to 0} \lambda^ {- 1} \psi^ {- 1} (\psi (h (\lambda x)) \psi (h (\lambda y))). \end{array}
\]

For \(|\lambda|\) sufficiently small, \(\lambda x \in V\) and \(\lambda y \in V\) and hence the above expression is equal to

\[
\begin{array}{r l} & \lim _ {\lambda \in \mathbf {K} ^ {*}, \lambda \to 0} \lambda^ {- 1} \psi^ {- 1} (f (\phi (\lambda x)) f (\phi (\lambda y))) \\ = & \lim _ {\lambda \in \mathbf {K} ^ {*}, \lambda \to 0} \lambda^ {- 1} (\psi^ {- 1} \circ f) (\phi (\lambda x) \phi (\lambda y)) \\ = & \lim _ {\lambda \in \mathbf {K} ^ {*}, \lambda \to 0} \lambda^ {- 1} g (\phi^ {- 1} (\phi (\lambda x) \phi (\lambda y))) \\ = & \lim _ {\lambda \in \mathbf {K} ^ {*}, \lambda \to 0} h (\lambda^ {- 1} \phi^ {- 1} (\phi (\lambda x) \phi (\lambda y))) \\ = & h (\lim _ {\lambda \in \mathbf {K} ^ {*}, \lambda \to 0} \lambda^ {- 1} \phi^ {- 1} (\phi (\lambda x) \phi (\lambda y))) \\ = & h (x + y). \end{array}
\]

Thus \(h\) is continuous and linear, hence \(g = h|\mathrm{V}\) is analytic, hence \(f\) is analytic on \(\phi(\mathrm{V})\) and hence \(f\) is analytic (§ 1, no. 10).

COROLLARY 1. Let G be a topological group. There exists on G at most one analytic manifold structure over R (resp. \(Q_{p}\) ) compatible with the topological group structure on G.

This follows immediately from Theorem 1.

DEFINITION 1. A topological group G is called a real (resp. p-adic) Lie group if there exists on G a real (resp. p-adic) Lie group structure compatible with its topology.

That structure is then unique and we can speak of the dimension of such a group. If G and H are two such groups, every continuous morphism of G into H is analytic.

COROLLARY 2. Let G be a topological group and V an open neighbourhood of e. Suppose that V has an analytic manifold structure which makes it into a real (resp. p-adic) Lie group germ. Then G is a real (resp. p-adic) Lie group.

Let \(g \in G\) . There exists an open neighbourhood \(V'\) of \(e\) in \(G\) such that \(V' \cup gV'g^{-1} \subset V\) . The mapping \(v \mapsto gvg^{-1}\) of \(V'\) into \(V\) is a continuous and hence analytic morphism of the Lie group germ \(V'\) into the Lie group germ \(V\) . It then suffices to apply Proposition 18 of § 1, no. 9.

Remarks. (1) Theorem 1 and its corollaries are no longer true if \(\mathbf{R}\) (resp. \(\mathbf{Q}_p\) ) is replaced by \(\mathbf{C}\) (Exercise 1).

\begin{enumerate}
\def\labelenumi{(\arabic{enumi})}
\setcounter{enumi}{1}
\tightlist
\item
  Let G be a topological group. It can be shown† that the following conditions are equivalent: (a) G is a finite-dimensional real Lie group; (b) G is locally compact and there exists a neighbourhood of e containing no subgroup distinct from \{e\}; (c) there exists an open neighbourhood of e homeomorphic to an open ball of a space \(R^{n}\) . (For a much less difficult result, cf.~Exercise 6.)
\end{enumerate}
† See for example D. Montgomery and L. Zippin, Topological transformation groups, Interscience Tracts in pure and applied mathematics, no. 1, Interscience publishers, New York 1955 (in particular pp. 169 and 184).

PROPOSITION 1. Let \(G, G'\) be topological groups and \(f\) a continuous morphism of \(G\) into \(G'\) . Suppose that one of the following three cases holds:

\begin{enumerate}
\def\labelenumi{(\alph{enumi})}
\item
  G is a real Lie group and \(\mathbf{G}'\) is a \(p\) -adic Lie group;
\item
  G is a p-adic Lie group and \(G'\) is a real Lie group;
\item
  G is a p-adic Lie group and \(G'\) is a \(p'\) -adic Lie group with \(p \neq p'\) .
\end{enumerate}

Then \(f\) is locally constant.

Case (a). Let \(G_0\) be the identity component of \(G\) . Then \(f(G_0)\) is a connected subgroup of \(G'\) and hence \(f(G_0) = \{e\}\) and \(G_0\) is open in \(G\) .

Case (b). Let \(V'\) be a neighbourhood of e in \(G'\) such that every subgroup of \(G'\) contained in \(V'\) reduces to \(\{e\}\) (§4, no. 2, Corollary 1 to Theorem 2). There exists a neighbourhood V of e in G such that \(f(V) \subset V'\) . Then there exists an open subgroup \(G_{1}\) of G such that \(G_{1} \subset V\) (§7, no. 1, Proposition 1). Then \(f(G_{1}) = \{e\}\) .

Case (c). By § 7, Theorem 4 and the Corollary to Proposition 8, there exists a neighbourhood \(\mathbf{V}'\) of \(e\) in \(\mathbf{G}'\) such that, for all \(x' \in \mathbf{V}' - \{e\}\) , \(x'^{p^n}\) does not tend to \(e\) as \(n\) tends to \(+\infty\) . There exists a neighbourhood \(V\) of \(e\) in \(G\) such that \(f(V) \subset V'\) . By § 7, Theorem 4 and Proposition 9, there exists an open subgroup \(G_1\) of \(G\) such that \(G_1 \subset V\) and such that, for all \(x \in G_1\) , \(x^{p^n}\) tends to \(e\) as \(n\) tends to \(+\infty\) . Then \(f(G_1) = \{e\}\) .

